%%%%%%%%%%%%%%%%%%%%%%%%%%%%%%%%%%%%%%%%%%%%%%%%%%%%%%%%%%%%%%%%%%%%%%
% Lớp 10 | Thống kê | Các số đặc trưng đo mức độ phân tán của mẫu số liệu
% Dạng 0D6?4-1: Câu hỏi lý thuyết
% Số câu: 3  (sinh tự động bởi tach_cau.py; từ nay thầy sửa trực tiếp file này)
%%%%%%%%%%%%%%%%%%%%%%%%%%%%%%%%%%%%%%%%%%%%%%%%%%%%%%%%%%%%%%%%%%%%%%

%%%%%%%%%%%%%%%% TRẮC NGHIỆM 4 PHƯƠNG ÁN %%%%%%%%%%%%%%%%

% nguồn: MaHoaNH_TracNghiem_Lop10
\begin{ex}%[0D6N4-1]%[ma:91b6a062]
	Để đánh giá mức độ phân tán của các số liệu thống kê so với số trung bình, ta dùng đại lượng nào sau đây? 
	\choice 
		{Số  trung bình}
		{Số trung vị} 
		 {Mốt}
		{\True Phương sai} 
	\loigiai{ 
		Dựa vào ý nghĩa của phương sai và độ lệch chuẩn để đo mức độ phân tán của các số liệu trong  mẫu quanh số trung bình.
	 }
\end{ex}

% nguồn: MaHoaNH_TracNghiem_Lop10
\begin{ex}%[0D6N4-1]%[ma:4b333ec3]
	Tìm phát biểu đúng về phương sai của một mẫu số  liệu.
	 \choice
		{Phương sai được sử dụng làm đại diện cho các số liệu của  mẫu}
		{\True Phương sai được sử dụng  để đánh giá mức độ phân tán của các số liệu thống kê (so với số trung bình)}
		{Phương sai được tính bằng tổng số phần tử  của một mẫu số liệu}
		{Phương sai là số liệu xuất hiện nhiều nhất (số liệu  có tần số lớn nhất) trong bảng các số
		liệu  thống kê}
	 \loigiai{
		Ta có ý nghĩa của phương sai là \lq\lq Phương sai được  sử dụng để đánh giá mức độ phân tán của các số liệu thống kê (so với số trung bình).\rq\rq
	 }
\end{ex}

% nguồn: MaHoaNH_TracNghiem_Lop10
\begin{ex}%[0D6N4-1]%[ma:5db221eb]
	Chọn  câu đúng trong các câu trả lời sau đây: Phương sai bằng
	 \choice
		{một nửa  của độ lệch chuẩn}
		{căn bậc hai của độ  lệch chuẩn}
		{hai  lần của độ lệch chuẩn}
		 {\True bình phương của độ lệch chuẩn}
	\loigiai{ 
		Ta  có
		\begin{itemize}
			\item Phương  sai là $s_x^2$.
			 \item Độ lệch chuẩn là $s_x=\sqrt{s_x^2}$.
		\end{itemize}
		Suy ra phương sai bằng bình  phương của độ lệch chuẩn.
	 }
\end{ex}
